\section{Significance of the study}

For waste handlers, meaning the staff of hotels and restaurants and the collection personnel, the prototype may make it easier to separate the supported types of waste and to see which compartments need attention. Current fill levels and estimated collection times may help collectors decide which sites to visit first, and the acknowledgement and collection records will let staff tell the difference between an alert that has been received and a bin that has actually been emptied. Any improvement in collection timing will have to be shown during testing.

For local authorities such as KCCA, records of waste categories, quantities and collection times from participating sites could fill part of the data gap that Okello \cite{okello} points out. KCCA has already asked for separation at source \cite{kccand}, and a site that separates its waste and records where it goes gives the authority something it can check. The study will also show whether this kind of monitoring is practical before anyone proposes it for more sites.

For the users of the system, which are the hotels and restaurants themselves, the study will show whether automatic sorting is practical in a place where items may be dirty, damaged or hard to recognise. It will also give them a record of what they throw away and who collects it, which they can use in their own planning.

For the environment, the benefit depends on how much of the waste is really separated and recovered. If food waste goes to farms or composting and recyclable material goes to recipients who accept it, less of it ends up in a landfill. The collapse at Kiteezi shows what happens when a disposal site is overloaded, so reducing what is sent there matters. The benefit will depend on what each recipient is willing to accept, and the study will look into those conditions before proposing any transfers.

For researchers and students, the project will produce a documented prototype and evaluation results that show both what it can do and where it falls short. Classification precision and recall, the proportion of items routed correctly, fill-level error, prediction error and alert delivery time will give measurable evidence. The findings may help guide later work on waste management systems that suit local operating conditions.
